\documentclass{article}%
\usepackage[T1]{fontenc}%
\usepackage[utf8]{inputenc}%
\usepackage{lmodern}%
\usepackage{textcomp}%
\usepackage{lastpage}%
\usepackage{geometry}%
\geometry{margin=2.2cm}%
%
\usepackage{amsmath}%
\usepackage{amssymb}%
\usepackage{booktabs}%
\usepackage{longtable}%
\usepackage{array}%
\usepackage{hyperref}%
\hypersetup{hidelinks}%
\title{ITA 2027 {-} Física\\[4pt]\large Comparação com o gabarito Poliedro Resolve (ITA 2027, 2ª fase, dia 3 {-} Física)\thanks{Gerado com gpt-resolve: \url{https://github.com/lgabs/gpt-resolve}}}%
\date{29/10/2026}%
%
\begin{document}%
\normalsize%
\maketitle%
\section*{Resumo}\textbf{Nota média: 10,0/10} (24 de 24 itens com a resolução correta; cada item vale o mesmo dentro da questão).\par\medskip%
\begin{longtable}{@{}c c p{0.78\textwidth}@{}}\toprule Questão & Nota & Resumo \\ \midrule
1 & 10,0 & A resolução e o gabarito do Poliedro concordam nos três itens. a) k\_c = MgL/2. b) São três regimes: k/k\_c $\geq$ 1 dá um equilíbrio estável; 2/$\pi$ < k/k\_c < 1 dá $\theta$=0 instável e $\pm$$\theta$* estáveis; k/k\_c $\leq$ 2/$\pi$ dá $\theta$=0 instável. c) Com U = k$\theta$$^2$/2 + (MgL/2) cos $\theta$, os gráficos são poço único, poço duplo e 'morro'. Refiz as contas e verifiquei numericamente a monotonicidade de U, o valor $\theta$* = 1,1311 para r = 0,8 e os valores de U nas bordas e nos mínimos; tudo está certo. A resolução é mais completa que o gabarito: trata o caso marginal k = k\_c, calcula U($\theta$*) e compara as bordas com MgL/2. Achei só um pequeno erro de digitação (parênteses faltando) na dedução do gabarito no item a, que não muda o resultado. \\
2 & 10,0 & A questão 2 está correta e bate com o gabarito do Poliedro nos dois itens. No item a, a resolução dá H\_n = 0,49\^{}n V$^2$/(8g), que é a mesma coisa que 0,7\^{}(2n) V$^2$/(8g) do gabarito. No item b, as duas chegam a $\Delta$t = V/{[}g(1-0,7){]} = 10V/(3g), somando a PG dos tempos de voo até a componente vertical se anular. Conferi as contas com Python e reli o enunciado no PDF original, que coincide com a transcrição. Não há erro de física nem de conta. As únicas ressalvas são de importância baixa: a palavra \textbackslash{}"energia\textbackslash{}" no enunciado é ambígua, chegar ao mínimo exige infinitos choques (uma idealização) e há erros de digitação na resolução do gabarito. \\
3 & 10,0 & Os dois itens da resolução são equivalentes ao gabarito do Poliedro, e os dois estão certos. Em (a) e (b), as expressões só diferem pela posição de R dentro das raízes: R$\cdot$$\omega$\_e = $\surd$R$\cdot$($\mu$e$^2$g$^2$ - $\alpha$$^2$R$^2$)\^{}(1/4). Conferi algebricamente, em 10$^{4}$ conjuntos de parâmetros aleatórios e por simulação direta, com D = 2,545 m e T = 2,366 s no caso teste. As únicas ressalvas são do gabarito: na resolução detalhada sumiram os quadrados de $\mu$\_e e g na linha (I) e na expressão final de D, e a forma fatorada final de (b) perdeu os parênteses. A resposta resumida do gabarito está correta, então isso não muda a nota. Não consegui gravar uma nota na memória da equipe: o sandbox desta sessão bloqueia a escrita fora da pasta de trabalho. \\
4 & 10,0 & Questão 4 correta. No item a, resolução e gabarito coincidem: h\_f = h\_i(1+M/m)\^{}\{3/5\}. No item b, a equação da resolução, m(h\_max-h\_i) = (3/2)(m+M)h\_i{[}1-(h\_i/h\_max)\^{}\{2/3\}{]}, é algebricamente equivalente à interpretação adiabática reversível do Poliedro. Conferi a equivalência simbolicamente e por simulação numérica (M = m dá h\_max $\approx$ 2,256h\_i). O único ponto em aberto é a ambiguidade do enunciado sobre a reversibilidade do gás no item b, que o próprio gabarito reconhece. A resolução declara explicitamente a hipótese que adota, e essa é a leitura mais natural de 'altura máxima'. \\
5 & 10,0 & A resolução da questão 5 concorda com o gabarito do Poliedro nos dois itens. No item a, as expressões de p' e de y' são idênticas, e as duas partes concluem que as imagens são virtuais e ficam no meio n1. No item b, a resolução dá $\Delta$y = ($\lambda$/2h)(D/n2 + p/n1), que com p ? D se reduz exatamente ao $\lambda$D/(2n2h) do gabarito. Conferi as duas respostas com contas independentes, incluindo um traçado exato pelo princípio de Fermat no item b, e ambas estão corretas. O único ponto ambíguo é se $\lambda$ é medido no vácuo ou no meio n1: as duas partes adotam o vácuo, e a resolução registra a alternativa. \\
6 & 10,0 & Os três itens batem com o gabarito do Poliedro, e os dois estão certos: a) q = 3C$\varepsilon$/4; b) compressão de 2d/3 (as placas ficam a d/3); c) k = 27C$\varepsilon$$^2$/(64d$^2$). Conferi o enunciado no PDF original (pág. 6). Refiz todas as contas e conferi numericamente. A conservação da carga dá C' = 3C, a força entre as placas é q$^2$/(2$\varepsilon$0A) e o balanço de energia fecha. A resolução ainda nota um ponto que o gabarito não comenta: com esse k, x = d/3 é exatamente o limite de estabilidade (W'' = 0). Isso não altera as respostas. \\
7 & 10,0 & A resolução e o gabarito coincidem nos dois itens. No a), |V\_o| = (V\_s/4)|G$\varepsilon$+$\alpha$$\Delta$T|. No b), são quatro extensômetros, com R2 e R3 na tração e R1 e R4 na compressão (ou o inverso), dando V\_o = V\_s G $\varepsilon$ sem efeito térmico. Conferi o circuito no PDF original e refiz as contas, de forma simbólica e numérica, incluindo a expressão exata V\_o = V\_sG$\varepsilon$/(1+$\alpha$$\Delta$T): está tudo correto. A resolução é mais completa que o gabarito, porque também justifica por que essa configuração é a ótima. Não encontrei erros. Ficam só observações menores sobre a interpretação do enunciado. \\
8 & 10,0 & A resolução está correta e coincide com o gabarito Poliedro nos três itens. a) v = $\surd$(2qV/m). b) Movimento helicoidal uniforme (MCU perpendicular a B mais MU ao longo de B), com o mesmo R e o mesmo T. c) cos$\theta$ = (BL/(2$\pi$n))$\surd$(q/(2mV)), n inteiro positivo, confirmado por integração numérica da trajetória. A única diferença é se o limite cos$\theta$ = 1 entra (resolução: estrito; gabarito: $\leq$), o que é irrelevante na prática. Não encontrei erro em nenhum dos dois. \\
9 & 10,0 & A resolução coincide com o gabarito do Poliedro nos três itens. a) E = E0{[}1+(-1)\^{}n{]} cos(kx-$\omega$t) $\hat z$, onda em +x. b) Para n par a interferência é construtiva (2E0) e para n ímpar é destrutiva total. c) B = -B0 cos(kx-$\omega$t) $\hat y$, com E$\times$B em +x?. Refiz as contas: a identidade em a) bate numericamente, a integração pela lei de Faraday e o produto vetorial $\hat z$$\times$(-$\hat y$)=x? estão certos, e não há divergência. A resolução está correta e é mais completa que o gabarito: verifica pela fórmula de soma em produto, pela lei de Faraday e pelo vetor de Poynting, e trata o caso n ímpar (B = 0). \\
10 & 10,0 & A resolução está correta nos dois itens e concorda com o gabarito do Poliedro. No item a, as duas dão a fórmula de Compton $\nu$' = $\nu$/{[}1 + (h$\nu$/(m\_e c$^2$))(1-cos$\theta$){]}. No item b, as duas dão o mesmo sentido: o observador vai ao encontro do fóton espalhado. As expressões de V têm formas diferentes, mas são algebricamente iguais, o que conferi numericamente; todas equivalem a V = c{[}(1+k)$^2$-1{]}/{[}(1+k)$^2$+1{]}, com k = h$\nu$(1-cos$\theta$)/(m\_e c$^2$). O gabarito tem dois erros de digitação na resolução detalhada (parênteses e quadrado fora do lugar), mas sua resposta final está correta. Não há divergência de mérito. \\
\bottomrule\end{longtable}%
\section*{Principais pontos de dúvida}%
\begin{itemize}%
\item%
\textbf{Q4} \textit{(statement)} No item b, o enunciado não diz se o gás evolui de forma reversível durante o movimento rápido do pistão, e o próprio Poliedro apresenta duas leituras. A resolução adota a isentrópica (pistão lento comparado com a velocidade do som), que coincide com a interpretação 1 do gabarito. Se a banca do ITA esperasse outra leitura, por exemplo a do repouso final h\_eq = (1+3M/5m)h\_i, a nota mudaria. A expressão 'altura máxima' e o pedido de 'uma equação que relaciona' (resposta implícita, transcendente) favorecem a leitura da resolução.%
\end{itemize}%
\section*{Questão 1}%
\href{https://poliedroresolve.sistemapoliedro.com.br/vestibulares/ita/2027/ita-2-fase-dia-3-22-10-2026-2-fase-dia-3/dissertativa-1-fisica-2-fase-dia-3-ita-2027}{Resolução do gabarito}\par\smallskip%
\begin{longtable}{@{}p{0.06\textwidth} p{0.33\textwidth} p{0.33\textwidth} p{0.1\textwidth} p{0.08\textwidth}@{}}\toprule Item & Resolução & Gabarito & Veredito & Certo \\ \midrule
a & k\_c = MgL/2 (o torque é $\tau$ $\approx$ -(k - MgL/2)$\theta$, então $\theta$=0 é estável se k > MgL/2; em k = k\_c continua estável pelo termo cúbico -k\_c$\theta$$^3$/6) & k\_c = MgL/2 & igual & ambos \\
b & k/k\_c $\geq$ 1: só $\theta$=0, estável. 2/$\pi$ < k/k\_c < 1: três equilíbrios, $\theta$=0 instável e $\pm$$\theta$* estáveis, com sen $\theta$*/$\theta$* = k/k\_c. 0 < k/k\_c $\leq$ 2/$\pi$: só $\theta$=0, instável (a haste tomba). & k/k\_c $\geq$ 1: um único equilíbrio, $\theta$=0, estável. 2/$\pi$ < k/k\_c < 1: três equilíbrios, $\theta$=0 instável e $\pm$$\theta$0 estáveis (k$\theta$0 = k\_c sen $\theta$0). k/k\_c $\leq$ 2/$\pi$: um único equilíbrio ($\theta$=0), instável, no intervalo aberto. & igual & ambos \\
c & U\_total = k$\theta$$^2$/2 + (MgL/2) cos $\theta$, com U(0) = MgL/2 e U($\pm$$\pi$/2) = k$\pi$$^2$/8. Esboços: (i) k $\geq$ k\_c: poço único com mínimo em 0; (ii) 2/$\pi$ < k/k\_c < 1: poço duplo em 'W', com máximo local em 0 e mínimos em $\pm$$\theta$*, onde U($\theta$*) = (MgL/2)(cos $\theta$* + ($\theta$*/2) sen $\theta$*) < MgL/2; (iii) k/k\_c $\leq$ 2/$\pi$: 'morro' (?) com máximo em 0, decrescente até as bordas. & U($\theta$) = (MgL/2) cos $\theta$ + k$\theta$$^2$/2. k $\geq$ k\_c: poço único com mínimo em 0; 2/$\pi$ < k/k\_c < 1: duplo poço (máximo local em 0, mínimos em $\pm$$\theta$0); k/k\_c $\leq$ 2/$\pi$: máximo em 0 e U decrescendo para as bordas. & igual & ambos \\
\bottomrule\end{longtable}%
\begin{itemize}%
\item%
\textbf{a)} Torque em relação ao mancal: $\tau$ = -k$\theta$ + (MgL/2) sen $\theta$ $\approx$ (MgL/2 - k)$\theta$, que é restaurador se k > MgL/2. Logo k\_c = MgL/2 e os dois chegam a isso. A resolução ainda confere pelo MHS (I\_O = ML$^2$/3) e analisa o caso limite k = k\_c, que continua estável por causa do termo de 4ª ordem de U. Isso é coerente com o enunciado, que diz que a estabilidade se perde para k < k\_c.%
\item%
\textbf{b)} Os equilíbrios saem de (k/k\_c)$\theta$ = sen $\theta$. Como sen $\theta$/$\theta$ é estritamente decrescente em (0, $\pi$/2), indo de 1 a 2/$\pi$, existe uma raiz não nula $\pm$$\theta$* só quando 2/$\pi$ < r < 1. As fronteiras batem nos dois: r = 1 entra no caso estável, e r = 2/$\pi$ entra no caso de equilíbrio único instável, porque a raiz $\pm$$\pi$/2 fica fora do intervalo aberto. Estabilidade: U'' = k\_c(r - cos $\theta$). Em $\theta$*, U'' = k\_c(sen $\theta$* - $\theta$* cos $\theta$*)/$\theta$* > 0, portanto $\theta$* é mínimo. Conferi numericamente: para r = 1,5 e r = 1, dU/d$\theta$ > 0 em (0, $\pi$/2); para r = 0,4 e r = 2/$\pi$, dU/d$\theta$ < 0 em (0, $\pi$/2). A condição sen $\theta$*/$\theta$* = k/k\_c da resolução é a mesma k$\theta$0 = k\_c sen $\theta$0 do gabarito.%
\item%
\textbf{c)} Os dois usam a mesma expressão de U e descrevem as mesmas três formas de gráfico. Conferi os detalhes extras da resolução. Para r = 0,8: $\theta$* = 1,1311 (bate com o tick do gráfico), U($\theta$*)/k\_c = 0,9374 < 1 e U($\pi$/2)/k\_c = 0,987. As bordas ficam acima de MgL/2 se r > 8/$\pi$$^2$ $\approx$ 0,8106. No caso (iii), U($\pi$/2) $\leq$ ($\pi$/4)$\cdot$MgL/2. As curvas do pgfplots (0,75x$^2$, 0,4x$^2$ e 0,2x$^2$, somadas a cos x) correspondem exatamente a U/k\_c para r = 1,5, 0,8 e 0,4. Em todos os casos a resolução está certa e é mais detalhada que o gabarito.%
\end{itemize}%
\paragraph{Pontos de dúvida}%
\begin{itemize}%
\item%
\textit{(key)} Erro de digitação na resolução do gabarito, item a: '$\tau$ $\approx$ MgL/2$\theta$ - k$\theta$ $\Rightarrow$ $\tau$ $\approx$ MgL/2 - k$\theta$'. Faltam parênteses: o certo é $\tau$ $\approx$ (MgL/2 - k)$\theta$. Isso não muda a conclusão k\_c = MgL/2.%
\item%
\textit{(statement)} O caso limite k = k\_c não é linearmente estável; é estável só pelo termo de 4ª ordem (U - U(0) $\approx$ k\_c$\theta$$^{4}$/24). Gabarito e resolução colocam esse caso junto com 'estável', o que é coerente com a frase do enunciado ('deixa de ser estável para k < k\_c'). Um corretor pode exigir ou ignorar essa justificativa.%
\item%
\textit{(statement)} Referência de altura: o enunciado manda tomar o 'mancal inferior' como zero, e a figura mostra o suporte/mancal abaixo do bloco de polímero. Os dois tomaram o zero no pivô (centro do polímero), o que dá U(0) = MgL/2. Se o zero fosse a base do suporte, U só ganharia uma constante aditiva, sem mudar a forma dos gráficos.%
\item%
\textit{(statement)} No caso k/k\_c = 2/$\pi$, as raízes $\pm$$\pi$/2 caem exatamente na borda do intervalo aberto. Os dois tratam esse caso como um único equilíbrio instável, o que está certo no intervalo pedido, mas é um ponto de fronteira que o corretor pode querer ver comentado. A resolução comenta.%
\end{itemize}%
\section*{Questão 2}%
\href{https://poliedroresolve.sistemapoliedro.com.br/vestibulares/ita/2027/ita-2-fase-dia-3-22-10-2026-2-fase-dia-3/dissertativa-2-fisica-2-fase-dia-3-ita-2027}{Resolução do gabarito}\par\smallskip%
\begin{longtable}{@{}p{0.06\textwidth} p{0.33\textwidth} p{0.33\textwidth} p{0.1\textwidth} p{0.08\textwidth}@{}}\toprule Item & Resolução & Gabarito & Veredito & Certo \\ \midrule
a & H\_n = (0,49)\^{}n V$^2$/(8g) & h\_n = 0,7\^{}(2n) V$^2$/(8g) & equivalente & ambos \\
b & $\Delta$t = V/{[}g(1-0,7){]} = 10V/(3g), com E\_min = 3mV$^2$/8 & T = (2V sen30$^\circ$/g)$\cdot$1/(1-e) = 10V/(3g) & igual & ambos \\
\bottomrule\end{longtable}%
\begin{itemize}%
\item%
\textbf{a)} 0,7\^{}(2n) = (0,7$^2$)\^{}n = 0,49\^{}n, então as duas expressões são a mesma. O raciocínio também é o mesmo: com o solo liso, v\_x se conserva; a cada choque v\_y é multiplicado por e, logo v\_y,n = e\^{}n$\cdot$V$\cdot$sen30$^\circ$ e H = v\_y$^2$/(2g). A conta confere.%
\item%
\textbf{b)} O valor é o mesmo e a interpretação também: a energia mínima é atingida depois de infinitos choques, quando a componente vertical se anula e só resta a horizontal. Os tempos de voo formam uma PG de razão e e primeiro termo 2V$\cdot$sen30$^\circ$/g = V/g, cuja soma é (V/g)/(1-0,7) = 10V/(3g). Conferi com frações exatas. A resolução ainda informa E\_min = 3mV$^2$/8 e x = 5$\surd$3V$^2$/(3g) nesse instante; os dois valores estão corretos.%
\end{itemize}%
\paragraph{Pontos de dúvida}%
\begin{itemize}%
\item%
\textit{(statement)} O enunciado fala em "mínimo da sua energia" sem dizer qual energia. A resolução e o gabarito entendem como energia mecânica total, que só cai nos choques e tende a 3mV$^2$/8. Também daria para ler como energia cinética ou potencial: a cinética é mínima no ápice do primeiro voo, em t = V/(2g), e a potencial é mínima a cada toque no solo. Essas leituras são menos naturais, mas um professor pode querer confirmar que a interpretação adotada é a esperada pela banca.%
\item%
\textit{(solution)} O mínimo só é atingido como limite de infinitos choques, cujos tempos formam uma série convergente (um 'paradoxo de Zenão' no modelo idealizado). A resolução e o gabarito tratam esse limite como um instante finito, T = 10V/(3g). Isso é padrão nesse tipo de questão, mas é uma idealização: e é tomado como constante para velocidades arbitrariamente pequenas e a partícula como pontual.%
\item%
\textit{(key)} A resolução do gabarito tem problemas de digitação na fórmula do item a: '(e\^{}nVsen$\theta$\^{}2)/2g' deveria ser (e\^{}n V sen$\theta$)$^2$/(2g), e há um 'gh' perdido em '...=(0,7\^{}2nV\^{}20,5\^{}2)/2gh = ...'. A resposta final, 0,7\^{}(2n)V$^2$/(8g), está correta.%
\end{itemize}%
\section*{Questão 3}%
\href{https://poliedroresolve.sistemapoliedro.com.br/vestibulares/ita/2027/ita-2-fase-dia-3-22-10-2026-2-fase-dia-3/dissertativa-3-fisica-2-fase-dia-3-ita-2027}{Resolução do gabarito}\par\smallskip%
\begin{longtable}{@{}p{0.06\textwidth} p{0.33\textwidth} p{0.33\textwidth} p{0.1\textwidth} p{0.08\textwidth}@{}}\toprule Item & Resolução & Gabarito & Veredito & Certo \\ \midrule
a & D = $\surd$(R$^2$ + (2hR/g)$\cdot$$\surd$($\mu$e$^2$g$^2$ - $\alpha$$^2$R$^2$)), com $\omega$\_s$^2$ = $\surd$($\mu$e$^2$g$^2$ - $\alpha$$^2$R$^2$)/R & D = R$\cdot$$\surd$(1 + (2h/g)$\cdot$$\surd$($\mu$e$^2$g$^2$/R$^2$ - $\alpha$$^2$)), com $\omega$\_e$^2$ = $\surd$($\mu$e$^2$g$^2$/R$^2$ - $\alpha$$^2$) & equivalente & ambos \\
b & T = $\surd$R$\cdot$($\mu$e$^2$g$^2$ - $\alpha$$^2$R$^2$)\^{}(1/4)$\cdot${[}1/($\alpha$R) + 1/($\mu$c g){]} + $\surd$(2h/g)  (t\_s + t\_q + t\_d, com $\mu$e g > $\alpha$R) & t\_p = $\omega$\_e/$\alpha$ + $\surd$(2h/g) + $\omega$\_e R/($\mu$c g), com $\omega$\_e = ($\mu$e$^2$g$^2$/R$^2$ - $\alpha$$^2$)\^{}(1/4) & equivalente & ambos \\
\bottomrule\end{longtable}%
\begin{itemize}%
\item%
\textbf{a)} É só tirar R de dentro da raiz: $\surd$($\mu$e$^2$g$^2$ - $\alpha$$^2$R$^2$)/R = $\surd$($\mu$e$^2$g$^2$/R$^2$ - $\alpha$$^2$), então $\omega$\_s$^2$ = $\omega$\_e$^2$ e (2hR/g)$\cdot$$\surd$($\mu$e$^2$g$^2$-$\alpha$$^2$R$^2$) = R$^2$$\cdot$(2h/g)$\cdot$$\surd$($\mu$e$^2$g$^2$/R$^2$-$\alpha$$^2$). Refiz o raciocínio. Enquanto não escorrega, o atrito estático fornece a aceleração m$\cdot$R$\surd$($\alpha$$^2$+$\omega$$^{4}$). O escorregamento ocorre quando essa força atinge $\mu$e$\cdot$m$\cdot$g. Daí vem um lançamento horizontal de altura h com v = $\omega$\_e$\cdot$R tangente à borda. No plano do solo, o deslocamento d = v$\surd$(2h/g) é perpendicular ao raio R, logo D = $\surd$(R$^2$+d$^2$). Conferi numericamente (g=9,8; R=2; $\alpha$=1,5; $\mu$e=0,6; $\mu$c=0,4; h=1,2): as duas fórmulas dão D = 2,545 m, o mesmo valor de uma simulação direta com a posição angular real $\theta$\_s = $\alpha$t\_s$^2$/2. Em 10$^{4}$ conjuntos de parâmetros aleatórios, a diferença relativa entre as duas formas ficou em torno de 10$^{-}$$^{1}$$^{5}$. O enunciado original (pág. 4 do PDF) bate com a transcrição, e a figura põe a origem no solo, sob o eixo do carrossel.%
\item%
\textbf{b)} $\surd$R$\cdot$($\mu$e$^2$g$^2$ - $\alpha$$^2$R$^2$)\^{}(1/4) = $\surd$R$\cdot$$\surd$R$\cdot$($\mu$e$^2$g$^2$/R$^2$ - $\alpha$$^2$)\^{}(1/4) = R$\cdot$$\omega$\_e. Assim, R$\cdot$$\omega$\_e/($\alpha$R) = $\omega$\_e/$\alpha$ = t\_e e R$\cdot$$\omega$\_e/($\mu$c g) = t\_f. As duas expressões são idênticas. As duas resoluções usam a mesma física: tempo até escorregar $\omega$\_e/$\alpha$, queda $\surd$(2h/g) e, no solo, só a componente vertical da velocidade é anulada no choque, seguida de MRUV com desaceleração $\mu$c$\cdot$g. No teste numérico, as duas dão T = 2,366 s, igual à simulação direta.%
\end{itemize}%
\paragraph{Pontos de dúvida}%
\begin{itemize}%
\item%
\textit{(key)} Na resolução do gabarito, a linha (I) traz $\omega$\_e$^2$ = $\surd$($\mu$\_e g/R - $\alpha$$^2$) e a expressão final de D traz $\surd$($\mu$\_e g/R$^2$ - $\alpha$$^2$): faltam os quadrados em $\mu$\_e e g, e as dimensões não fecham. A resposta resumida do gabarito está correta ($\mu$e$^2$g$^2$/R$^2$ - $\alpha$$^2$). Parece expoente perdido na transcrição, não erro de física, mas vale conferir no original do Poliedro.%
\item%
\textit{(key)} A forma fatorada final do gabarito em (b), 't\_P = 1/$\alpha$ + R/$\mu$\_cg root{[}4{]}(...) + $\surd$(2h/g)', saiu com os parênteses perdidos. A leitura pretendida é (1/$\alpha$ + R/($\mu$c g))$\cdot$$^{4}$$\surd$($\mu$e$^2$g$^2$/R$^2$ - $\alpha$$^2$) + $\surd$(2h/g), que bate com a soma t\_e + t\_q + t\_f escrita logo antes.%
\item%
\textit{(statement)} Hipótese no choque com o solo: as duas resoluções supõem que só a componente vertical da velocidade é anulada, sem impulso de atrito durante o impacto. Essa é a leitura natural de 'não quica, mas realiza um deslizamento'. Com atrito de Coulomb no impacto, ficaria T = t\_s + v/($\mu$c g), como a própria resolução observa. Seria uma resposta diferente, mas improvável de ser a esperada pela banca.%
\item%
\textit{(statement)} Idealizações comuns às duas resoluções e não ditas no enunciado: a massa é pontual, sai do disco no instante em que o atrito estático satura (por estar na borda) com a velocidade tangencial da borda, e a espessura do disco e a resistência do ar são desprezadas. Também é preciso que $\mu$e$\cdot$g > $\alpha$$\cdot$R. A resolução declara essa condição, o gabarito não.%
\end{itemize}%
\section*{Questão 4}%
\href{https://poliedroresolve.sistemapoliedro.com.br/vestibulares/ita/2027/ita-2-fase-dia-3-22-10-2026-2-fase-dia-3/dissertativa-4-fisica-2-fase-dia-3-ita-2027}{Resolução do gabarito}\par\smallskip%
\begin{longtable}{@{}p{0.06\textwidth} p{0.33\textwidth} p{0.33\textwidth} p{0.1\textwidth} p{0.08\textwidth}@{}}\toprule Item & Resolução & Gabarito & Veredito & Certo \\ \midrule
a & h\_f = h\_i$\cdot$((m+M)/m)\^{}\{3/5\} (adiabática reversível, $\gamma$ = 5/3) & h\_f = (1 + M/m)\^{}\{3/5\}$\cdot$h\_i (adiabática reversível, $\gamma$ = 5/3) & igual & ambos \\
b & m(h\_max - h\_i) = (3/2)(m+M)h\_i{[}1 - (h\_i/h\_max)\^{}\{2/3\}{]}. Na forma polinomial, com y = (h\_max/h\_i)\^{}\{1/3\}: 2y$^{4}$+2y$^3$+2y$^2$ = 3(1+M/m)(y+1), com raiz única y > 1. Para M = m, h\_max $\approx$ 2,26 h\_i. & Forma geral (1ª lei): h\_max = (3M+5m)g h\_i/(3p\_fA + 2mg), com p\_f indeterminado. Supondo adiabática reversível: 3(M+m)(h\_i/h\_max)\^{}\{5/3\}h\_max + 2m h\_max = (3M+5m)h\_i. Alternativa (repouso final após amortecimento): h\_eq = (1 + 3M/(5m))h\_i. & equivalente & ambos \\
\bottomrule\end{longtable}%
\begin{itemize}%
\item%
\textbf{a)} Mesma resposta, já que (m+M)/m = 1 + M/m. Usa{-}se p\_iV\_i\^{}$\gamma$ = p\_fV\_f\^{}$\gamma$ com p\_i = (m+M)g/A, p\_f = mg/A e $\gamma$ = 5/3. O enunciado (pág. 5 do PDF) confere com a transcrição.%
\item%
\textbf{b)} A equação da resolução é algebricamente idêntica à interpretação 1 do Poliedro. Multiplicando por 2 e rearranjando, vem 2m$\cdot$h\_max + 3(m+M)h\_i\^{}\{5/3\}h\_max\^{}\{{-}2/3\} = (3M+5m)h\_i, que é a equação do gabarito. Testei com valores aleatórios e a diferença deu zero. Conferi também a parte numérica. Para M = m, a raiz da quártica dá h\_max/h\_i = 2,2559, e uma integração RK4 da equação de movimento do pistão, m$\cdot$? = p\_i(h\_i/h)\^{}\{5/3\}A - mg, dá o mesmo primeiro ponto de retorno. A redução a 2y$^{4}$+2y$^3$+2y$^2$ - 3k(y+1) = 0, a convexidade (f'' > 0, f(1) = -6M/m < 0) e o limite M ? m (h\_max $\approx$ h\_i(1 + 1,2M/m), contra 1,012 exato para M/m = 0,01) estão corretos. A forma geral do gabarito é válida, mas deixa p\_f solto, e o próprio gabarito admite isso. A alternativa 2 (h\_eq = 1,6h\_i para M = m) é a altura de repouso final após dissipação, não a altura máxima. Como o item pede 'a equação que relaciona h\_max e m, M e h\_i', a leitura isentrópica, adotada pela resolução, é a mais natural.%
\end{itemize}%
\paragraph{Pontos de dúvida}%
\begin{itemize}%
\item%
\textit{(statement)} No item b, o enunciado não diz se o gás evolui de forma reversível durante o movimento rápido do pistão, e o próprio Poliedro apresenta duas leituras. A resolução adota a isentrópica (pistão lento comparado com a velocidade do som), que coincide com a interpretação 1 do gabarito. Se a banca do ITA esperasse outra leitura, por exemplo a do repouso final h\_eq = (1+3M/5m)h\_i, a nota mudaria. A expressão 'altura máxima' e o pedido de 'uma equação que relaciona' (resposta implícita, transcendente) favorecem a leitura da resolução.%
\item%
\textit{(key)} A resposta principal do gabarito no item b, h\_max = (3M+5m)g h\_i/(3p\_fA+2mg), depende de p\_f, que não está entre as variáveis pedidas (m, M, h\_i). Só fica completa com uma hipótese adicional. A alternativa 2 (h\_eq) responde outra pergunta: altura de equilíbrio final, não altura máxima.%
\item%
\textit{(key)} A transcrição do gabarito traz 'p\_i = (M + mg)/A' e '(3M + 5mg)', em vez de (M+m)g/A e (3M+5m)g. São erros de extração ou digitação sem efeito no resultado.%
\item%
\textit{(solution)} A estimativa h\_max $\approx$ (3/2)(M/m)h\_i para M ? m vale só em primeira ordem. Para M = 100m, o valor exato é $\approx$ 147h\_i, contra 150h\_i da aproximação, ou 151,5h\_i se for usado 1,5(1+M/m). Isso não afeta a resposta final.%
\end{itemize}%
\section*{Questão 5}%
\href{https://poliedroresolve.sistemapoliedro.com.br/vestibulares/ita/2027/ita-2-fase-dia-3-22-10-2026-2-fase-dia-3/dissertativa-5-fisica-2-fase-dia-3-ita-2027}{Resolução do gabarito}\par\smallskip%
\begin{longtable}{@{}p{0.06\textwidth} p{0.33\textwidth} p{0.33\textwidth} p{0.1\textwidth} p{0.08\textwidth}@{}}\toprule Item & Resolução & Gabarito & Veredito & Certo \\ \midrule
a & p' = -n2 p R/(n1 R + (n2 - n1) p): imagens virtuais no meio n1, a |p'| = n2pR/(n1R+(n2-n1)p) à esquerda do vértice, entre as fontes e C; y' = $\pm$ n1 R h/(n1 R + (n2 - n1) p), imagens direitas e reduzidas & p' = -n2 p R/(n1 R + (n2 - n1) p) (imagens virtuais no meio n1); y' = $\pm$ n1 R h/(n1 R + (n2 - n1) p) & igual & ambos \\
b & $\Delta$y = ($\lambda$/2h)(D/n2 + p/n1) $\approx$ $\lambda$D/(2 n2 h), com $\lambda$ no vácuo & y = $\lambda$D/(2 n2 h) & equivalente & ambos \\
\bottomrule\end{longtable}%
\begin{itemize}%
\item%
\textbf{a)} As duas usam a equação do dioptro com raio algébrico -R, porque C está no meio incidente, e chegam às mesmas expressões. Conferi numericamente que p' satisfaz n1/p + n2/p' = (n2-n1)/(-R) e que A = -n1p'/(n2p) = n1R/(n1R+(n2-n1)p). A resolução ainda deduz y' pelo raio que passa por C e confirma que as imagens ficam entre as fontes e C. As contas extras estão corretas: |p'|-p = p(n2-n1)(R-p)/(\ldots{}) > 0 e R-|p'| = n1R(R-p)/(\ldots{}) > 0.%
\item%
\textbf{b)} A forma aproximada coincide com o gabarito. A forma completa da resolução, com a distância efetiva L = D + n2p/n1 até as fontes virtuais, é mais precisa: o termo $\lambda$p/(2n1h) desaparece quando p ? D, como o enunciado manda considerar. Conferi com Fermat exato para dioptro plano (n1=1,2; n2=1,7; p=5 cm; D=2 m; h=0,4 mm; $\lambda$=600 nm). O traçado exato dá 0,913603 mm, a forma completa dá 0,913603 mm e $\lambda$D/(2n2h) dá 0,882 mm. As duas respostas valem dentro da aproximação pedida.%
\end{itemize}%
\paragraph{Pontos de dúvida}%
\begin{itemize}%
\item%
\textit{(statement)} O enunciado não diz se $\lambda$ é o comprimento de onda no vácuo ou no meio n1, onde estão as fontes. Tanto o gabarito quanto a resolução tomam $\lambda$ no vácuo. Se $\lambda$ for o comprimento no meio n1, a resposta vira n1$\lambda$D/(2n2h). A resolução cita essa alternativa explicitamente.%
\item%
\textit{(solution)} No item b, a resolução dá primeiro a forma completa ($\lambda$/2h)(D/n2 + p/n1) e só depois aproxima para $\lambda$D/(2n2h). As duas formas estão certas, e a final coincide com o gabarito. Fica só o registro de que a resposta de referência é a aproximada, com p ? D.%
\item%
\textit{(key)} No item b, o argumento do gabarito usa raios paralelos (campo distante) e trata como iguais os trechos dos dois raios no meio n1. Isso só é válido no limite p ? D, o que o gabarito assume corretamente. A correção exata é o termo p/n1, que o gabarito descarta de forma implícita.%
\end{itemize}%
\section*{Questão 6}%
\href{https://poliedroresolve.sistemapoliedro.com.br/vestibulares/ita/2027/ita-2-fase-dia-3-22-10-2026-2-fase-dia-3/dissertativa-6-fisica-2-fase-dia-3-ita-2027}{Resolução do gabarito}\par\smallskip%
\begin{longtable}{@{}p{0.06\textwidth} p{0.33\textwidth} p{0.33\textwidth} p{0.1\textwidth} p{0.08\textwidth}@{}}\toprule Item & Resolução & Gabarito & Veredito & Certo \\ \midrule
a & q = 3C$\varepsilon$/4 (C' = 3C, V3 = $\varepsilon$/4; +q numa placa, -q na outra, C$\varepsilon$/4 fica no capacitor) & q = 3C$\varepsilon$/4 (C' = 3C, V3 = $\varepsilon$/4) & igual & ambos \\
b & Compressão $\Delta$x = 2d/3 (distância final d/3) & Compressão $\Delta$x = 2d/3 (distância final d/3) & igual & ambos \\
c & k = 27C$\varepsilon$$^2$/(64d$^2$) & k = 27C$\varepsilon$$^2$/(64d$^2$) & igual & ambos \\
\bottomrule\end{longtable}%
\begin{itemize}%
\item%
\textbf{a)} Conservação da carga: C$\varepsilon$ = (C + C')$\cdot$$\varepsilon$/4 $\Rightarrow$ C' = 3C, logo q = C'$\cdot$V3 = 3C$\varepsilon$/4. Conferi também que q + C$\varepsilon$/4 = C$\varepsilon$. A resolução e o gabarito seguem o mesmo raciocínio.%
\item%
\textbf{b)} C' = $\varepsilon$0A/x = Cd/x = 3C $\Rightarrow$ x = d/3. Como a mola começa no comprimento natural, $\Delta$x = d - d/3 = 2d/3, e é compressão porque as placas têm cargas opostas e se atraem. A resolução ainda observa que a deformação total não depende de como as placas se deslocam.%
\item%
\textbf{c)} Cada placa sente só o campo da outra, $\sigma$/(2$\varepsilon$0). Daí F = q$^2$/(2$\varepsilon$0A) = (9C$^2$$\varepsilon$$^2$/16)/(2Cd) = 9C$\varepsilon$$^2$/(32d). No equilíbrio, k$\cdot$(2d/3) = 9C$\varepsilon$$^2$/(32d) $\Rightarrow$ k = 27C$\varepsilon$$^2$/(64d$^2$). Conferi numericamente com C = $\varepsilon$ = d = 1: k = 0,421875 = 27/64. O balanço de energia da resolução também está certo: a energia inicial é 16/32$\cdot$C$\varepsilon$$^2$, a final é 4/32 + 3/32 = 7/32$\cdot$C$\varepsilon$$^2$, e a diferença é dissipada.%
\end{itemize}%
\paragraph{Pontos de dúvida}%
\begin{itemize}%
\item%
\textit{(statement)} Com k = 27C$\varepsilon$$^2$/(64d$^2$), a energia W(x) = Q0$^2$x/{[}2C(x+d){]} + $\frac12$k(d-x)$^2$ tem W' = 0 e W'' = 0 em x = d/3, com W''' $\neq$ 0. Conferi numericamente. O equilíbrio dado é, portanto, um ponto de inflexão: semiestável, exatamente no limite de pull{-}in. Em geral, W'' = 0 no equilíbrio só quando x = d/3, ou seja, quando V3 = $\varepsilon$/4. Para V3 < $\varepsilon$/4 o equilíbrio seria instável. Parece que a banca escolheu o caso{-}limite, talvez de propósito. Na prática, um sistema com inércia passaria do ponto e as placas colapsariam, e um sistema superamortecido só chegaria a d/3 assintoticamente. As respostas não mudam, porque saem das condições de equilíbrio dadas. A resolução comenta isso corretamente, e o gabarito não menciona.%
\item%
\textit{(statement)} Hipóteses implícitas que a resolução e o gabarito adotam: voltímetro, fios e chaves ideais; nenhum efeito de borda (d ? a, e depois d/3 ? a); efeito pendular dos fios desprezado, já que o enunciado não dá massas nem comprimentos; e, na direção horizontal, só a força elétrica e a força da mola.%
\item%
\textit{(key)} A resolução do gabarito veio com a formatação corrompida na extração. Por exemplo, aparece 'C$\varepsilon$ = C + C'$\varepsilon$/4C' = 3C' no lugar de C$\varepsilon$ = (C + C')$\varepsilon$/4 $\Rightarrow$ C' = 3C, e 'k2d/3 = 1/2$\varepsilon$\_0A3C$\varepsilon$/4\^{}2' no lugar de k$\cdot$2d/3 = (3C$\varepsilon$/4)$^2$/(2$\varepsilon$0A). As contas por trás estão certas. É só problema de transcrição.%
\end{itemize}%
\section*{Questão 7}%
\href{https://poliedroresolve.sistemapoliedro.com.br/vestibulares/ita/2027/ita-2-fase-dia-3-22-10-2026-2-fase-dia-3/dissertativa-7-fisica-2-fase-dia-3-ita-2027}{Resolução do gabarito}\par\smallskip%
\begin{longtable}{@{}p{0.06\textwidth} p{0.33\textwidth} p{0.33\textwidth} p{0.1\textwidth} p{0.08\textwidth}@{}}\toprule Item & Resolução & Gabarito & Veredito & Certo \\ \midrule
a & |V\_o| = (V\_s/4)$\cdot$|G$\varepsilon$ + $\alpha$$\Delta$T|, qualquer que seja o braço do extensômetro (sinal negativo em R1 ou R4, positivo em R2 ou R3) & |V0| = (Vs/4)$\cdot$|G$\varepsilon$ + $\alpha$$\Delta$T| & igual & ambos \\
b & Quatro extensômetros, nenhum resistor fixo: R2 e R3 em tração, R1 e R4 em compressão (ou o inverso). Assim V\_o = V\_s G $\varepsilon$, sem efeito térmico. & Quatro extensômetros: R2 e R3 na tração, R1 e R4 na compressão (ou tudo trocado), dando V0 = Vs G $\varepsilon$, sem efeito térmico & igual & ambos \\
\bottomrule\end{longtable}%
\begin{itemize}%
\item%
\textbf{a)} Pela figura conferida no PDF, R1 está entre D e A, R3 entre D e C, R2 entre A e B, R4 entre B e C, e a fonte está entre A e C. Então V\_D = V\_s R3/(R1+R3) e V\_B = V\_s R4/(R2+R4). Linearizando, V\_o $\approx$ (V\_s/4)(x2+x3-x1-x4), igual na resolução e no gabarito. Conferi numericamente com a expressão exata (G=2, $\varepsilon$=1e{-}5, $\alpha$$\Delta$T=3e{-}5): o extensômetro em cada um dos quatro braços dá |V\_o| $\approx$ (V\_s/4)(G$\varepsilon$+$\alpha$$\Delta$T), com sinal - em R1/R4 e + em R2/R3. A resolução ainda explica a dependência de sinal com a posição, o que é correto e vai além do que o gabarito pede.%
\item%
\textbf{b)} A configuração e o resultado são os mesmos. Conferência exata: com R1=R4=R0(1-G$\varepsilon$+$\alpha$$\Delta$T) e R2=R3=R0(1+G$\varepsilon$+$\alpha$$\Delta$T), sai V\_o = V\_s G$\varepsilon$/(1+$\alpha$$\Delta$T) $\approx$ V\_s G$\varepsilon$. Bate com o número calculado (1,99994e{-}5 contra 2e{-}5 para V\_s=1). A resolução ainda justifica por que essa é a configuração ótima: o módulo máximo 4 do coeficiente de G$\varepsilon$ só é atingido com n\_i=1 e sinais alternados, e o termo térmico se anula só se n2+n3=n1+n4. O gabarito apenas mostra que a configuração funciona.%
\end{itemize}%
\paragraph{Pontos de dúvida}%
\begin{itemize}%
\item%
\textit{(statement)} No item b, a expressão 'todo o sistema está na mesma temperatura' pode ser lida como $\Delta$T igual em todos os dispositivos (leitura da resolução e do gabarito) ou como $\Delta$T = 0. A resposta (ponte completa com 4 extensômetros) não muda em nenhuma das duas leituras, mas na segunda a justificativa da compensação térmica perde o sentido.%
\item%
\textit{(solution)} A polaridade da fonte não é dada explicitamente. Pela figura, o traço mais longo da bateria está em cima, o que indica V\_A > V\_C. A resolução assume isso e observa que inverter a fonte só troca o sinal de V\_o. Isso não afeta o item a (pede o módulo) nem a escolha do item b.%
\item%
\textit{(statement)} No item b, uma meia ponte (dois extensômetros em braços adjacentes, um na tração e outro na compressão) também cancelaria o efeito térmico, com sensibilidade 2 vezes a do item a. O critério 'máxima precisão' leva à ponte completa, que é a resposta da resolução e do gabarito. Um corretor poderia discutir crédito parcial para a meia ponte.%
\end{itemize}%
\section*{Questão 8}%
\href{https://poliedroresolve.sistemapoliedro.com.br/vestibulares/ita/2027/ita-2-fase-dia-3-22-10-2026-2-fase-dia-3/dissertativa-8-fisica-2-fase-dia-3-ita-2027}{Resolução do gabarito}\par\smallskip%
\begin{longtable}{@{}p{0.06\textwidth} p{0.33\textwidth} p{0.33\textwidth} p{0.1\textwidth} p{0.08\textwidth}@{}}\toprule Item & Resolução & Gabarito & Veredito & Certo \\ \midrule
a & v = $\surd$(2qV/m) & v = $\surd$(2qV/m) & igual & ambos \\
b & Movimento helicoidal uniforme: MRU ao longo de B com v cos$\theta$, mais MCU no plano perpendicular a B, com R = mv sen$\theta$/(qB) e T = 2$\pi$m/(qB). A resolução também dá o passo p = (2$\pi$ cos$\theta$/B)$\surd$(2mV/q). & Movimento helicoidal uniforme: MCU no plano perpendicular a B (R = mv sen$\theta$/(qB), T = 2$\pi$m/(qB)) composto com MU ao longo de B (v cos$\theta$). & igual & ambos \\
c & cos$\theta$ = BL/(2$\pi$n)$\cdot$$\surd$(q/(2mV)), com n = 1, 2, 3, \ldots{} e n > (BL/2$\pi$)$\surd$(q/(2mV)), para que cos$\theta$ < 1. & cos$\theta$ = qBL/(2$\pi$nmv) = (BL/(2$\pi$n))$\cdot$$\surd$(q/(2mV)), com n $\in$ $\mathbb{N}$* e cos$\theta$ $\leq$ 1, ou seja, n $\geq$ (BL/2$\pi$)$\surd$(q/(2mV)). & igual & ambos \\
\bottomrule\end{longtable}%
\begin{itemize}%
\item%
\textbf{a)} Pelo teorema da energia cinética, partindo do repouso, qV = mv$^2$/2. A resolução e o gabarito chegam à mesma expressão.%
\item%
\textbf{b)} A conclusão e a justificativa são as mesmas: a força magnética é perpendicular a B, então não há aceleração ao longo de B; e é perpendicular a v, então não realiza trabalho e faz o papel de resultante centrípeta. R e T coincidem. O passo e os casos{-}limite $\theta$ = 90$^\circ$ e $\theta$ = 0 são acréscimos corretos da resolução.%
\item%
\textbf{c)} As duas usam L = n$\cdot$v cos$\theta$$\cdot$T, com T = 2$\pi$m/(qB), e chegam à mesma expressão. Integrei a trajetória numericamente (RK4, q = 1,6e{-}19 C, m = 1,67e{-}27 kg, V = 1000 V, B = 0,1 T, L = 0,5 m, n = 3): depois de 3T a partícula está em (0,5; \textasciitilde{}0; \textasciitilde{}0), isto é, exatamente no ponto 2. A única diferença é a restrição complementar: a resolução exige desigualdade estrita (cos$\theta$ < 1, porque a partícula de fato sai da linha), e o gabarito aceita a igualdade (cos$\theta$ $\leq$ 1). Com cos$\theta$ = 1 ($\theta$ = 0), a partícula nunca sai da linha de campo, então nem há 'retorno'. A figura mostra $\theta$ $\neq$ 0. A diferença é irrelevante para a resposta pedida.%
\end{itemize}%
\paragraph{Pontos de dúvida}%
\begin{itemize}%
\item%
\textit{(solution)} No item c, a resolução restringe n com desigualdade estrita (n > (BL/2$\pi$)$\surd$(q/(2mV)), ou seja, cos$\theta$ < 1), enquanto o gabarito usa n $\geq$ (cos$\theta$ $\leq$ 1). O caso{-}limite cos$\theta$ = 1 é $\theta$ = 0: a partícula se move sobre a própria linha de campo, passa por 2 de qualquer forma e não há 'retorno'. As duas formas são defensáveis e não devem afetar a nota.%
\item%
\textit{(statement)} O enunciado não diz se o peso deve ser desprezado nem se o ponto 2 está no sentido de B. A resolução declara as duas hipóteses: peso desprezível e ponto 2 no sentido de B, pois cos$\theta$ > 0 pela figura. O gabarito as adota implicitamente.%
\end{itemize}%
\section*{Questão 9}%
\href{https://poliedroresolve.sistemapoliedro.com.br/vestibulares/ita/2027/ita-2-fase-dia-3-22-10-2026-2-fase-dia-3/dissertativa-9-fisica-2-fase-dia-3-ita-2027}{Resolução do gabarito}\par\smallskip%
\begin{longtable}{@{}p{0.06\textwidth} p{0.33\textwidth} p{0.33\textwidth} p{0.1\textwidth} p{0.08\textwidth}@{}}\toprule Item & Resolução & Gabarito & Veredito & Certo \\ \midrule
a & E = E0{[}1+(-1)\^{}n{]} cos(kx - $\omega$t) $\hat z$; é uma onda progressiva no sentido +x (velocidade $\omega$/k = c) com amplitude E0{[}1+(-1)\^{}n{]} & E = E0{[}1 + (-1)\^{}n{]} cos(kx - $\omega$t) $\hat z$; amplitude E0{[}1 + (-1)\^{}n{]}, onda no sentido +x & igual & ambos \\
b & $\Delta$$\varphi$ = n$\pi$ (d = n$\lambda$/2). n par: interferência construtiva (amplitude 2E0, intensidade 4 vezes a de cada onda). n ímpar: interferência destrutiva total (E = 0) & n par: construtiva (amplitude 2E0); n ímpar: destrutiva total (amplitude 0) & igual & ambos \\
c & B = -B0 cos(kx - $\omega$t) $\hat y$, com B0 = E0{[}1+(-1)\^{}n{]}/c (2E0/c para n par; B = 0 para n ímpar). B fica ao longo do eixo y, em fase com E, e aponta para -$\hat y$ quando E aponta para +$\hat z$, de modo que E$\times$B fica em +x? & B = -B0 cos(kx - $\omega$t) $\hat y$ (E $\times$ B no sentido +x?) & igual & ambos \\
\bottomrule\end{longtable}%
\begin{itemize}%
\item%
\textbf{a)} Mesma expressão e mesma justificativa: cos($\varphi$+n$\pi$) = (-1)\^{}n cos $\varphi$. Conferi numericamente para n = 0\ldots{}5 (erro \textasciitilde{}1e{-}15). A resolução também confere pela fórmula de soma em produto e justifica o sentido de propagação pela fase constante (x = x0 + ($\omega$/k)t). Está correta e mais completa que o gabarito.%
\item%
\textbf{b)} Mesma conclusão. Os acréscimos da resolução (d = n$\lambda$/2, I $\propto$ E$^2$ e portanto 4I0) estão corretos e não contradizem o gabarito.%
\item%
\textbf{c)} Mesma expressão vetorial, sentido e sinal. $\hat z$ $\times$ (-$\hat y$) = +x? está correto. A resolução confirma pela lei de Faraday: $\partial$B\_y/$\partial$t = $\partial$E\_z/$\partial$x leva a B\_y = -(E\_R/c) cos $\varphi$. Também acrescenta que B0 = E\_R/c, o que é coerente com o enunciado ('amplitude de módulo B0' do campo resultante). Ela ainda dá a forma alternativa caso B0 seja a amplitude de cada onda individual. Isso não altera a resposta principal.%
\end{itemize}%
\paragraph{Pontos de dúvida}%
\begin{itemize}%
\item%
\textit{(statement)} No item c), quando n é ímpar o campo resultante é nulo (B = 0). Nesse caso não existe 'orientação' de B, e a 'amplitude B0' é zero. O enunciado (e o gabarito) supõem tacitamente que B0 $\neq$ 0, ou seja, n par. A resolução trata os dois casos explicitamente, o que é adequado.%
\item%
\textit{(solution)} A resolução escreve B0 em função de E0 (B0 = E0{[}1+(-1)\^{}n{]}/c), enquanto o enunciado dá B0 como dado. Isso não é erro, porque a relação B0 = E\_R/c é física correta e a resposta pedida (-B0 cos(kx-$\omega$t) $\hat y$) aparece literalmente. Só vale notar que a resolução vai além do pedido e ainda dá uma forma alternativa para o caso de B0 ser a amplitude de cada onda individual.%
\item%
\textit{(key)} O gabarito justifica o sentido de propagação só dizendo que cos(kx-$\omega$t) = cos($\omega$t-kx) 'representa uma onda que se propaga no sentido +x', sem argumento de fase constante. A justificativa é suficiente para o nível da prova, mas é mais sucinta que a da resolução.%
\end{itemize}%
\section*{Questão 10}%
\href{https://poliedroresolve.sistemapoliedro.com.br/vestibulares/ita/2027/ita-2-fase-dia-3-22-10-2026-2-fase-dia-3/dissertativa-10-fisica-2-fase-dia-3-ita-2027}{Resolução do gabarito}\par\smallskip%
\begin{longtable}{@{}p{0.06\textwidth} p{0.33\textwidth} p{0.33\textwidth} p{0.1\textwidth} p{0.08\textwidth}@{}}\toprule Item & Resolução & Gabarito & Veredito & Certo \\ \midrule
a & $\nu$' = $\nu$ / {[}1 + (h$\nu$/(m\_e c$^2$))(1 - cos $\theta$){]} (Compton), obtida pela conservação da energia relativística e do momento, com a lei dos cossenos. & $\nu$' = $\nu$/(1 + (h$\nu$/(m\_e c$^2$))(1 - cos$\theta$)) (Compton) & igual & ambos \\
b & V = c($\nu$$^2$ - $\nu$'$^2$)/($\nu$$^2$ + $\nu$'$^2$) = c$\cdot$k(2+k)/(2+2k+k$^2$), com k = h$\nu$(1-cos$\theta$)/(m\_e c$^2$). Explicitamente, V = c\{{[}m\_e c$^2$ + h$\nu$(1-cos$\theta$){]}$^2$ - (m\_e c$^2$)$^2$\}/\{{[}m\_e c$^2$ + h$\nu$(1-cos$\theta$){]}$^2$ + (m\_e c$^2$)$^2$\}. O observador se move na direção do fóton espalhado, em sentido oposto ao de sua propagação (indo ao encontro dele). & O observador se move na direção do fóton espalhado, em sentido oposto ao de sua propagação (aproximando{-}se), com $\beta$ = ($\nu$$^2$ - $\nu$'$^2$)/($\nu$$^2$ + $\nu$'$^2$), ou seja, V = c$\cdot${[}(1 + h$\nu$(1-cos$\theta$)/(m\_e c$^2$))$^2$ - 1{]}/{[}(1 + h$\nu$(1-cos$\theta$)/(m\_e c$^2$))$^2$ + 1{]} & equivalente & ambos \\
\bottomrule\end{longtable}%
\begin{itemize}%
\item%
\textbf{a)} As duas fórmulas são idênticas e as deduções seguem o mesmo caminho: conservação da energia, triângulo de momentos com a lei dos cossenos e E$^2$ = p$^2$c$^2$ + m$^2$c$^{4}$, chegando a m\_e c$^2$($\nu$ - $\nu$') = h$\nu$$\nu$'(1 - cos $\theta$). A resolução ainda confere os casos limite ($\theta$ = 0 e h$\nu$ ? m\_e c$^2$). Conferi a transcrição contra o PDF original e ela está fiel.%
\item%
\textbf{b)} O sentido é o mesmo nas duas: ao encontro do fóton, porque $\nu$ > $\nu$' exige desvio Doppler para o azul. As expressões de V são algebricamente iguais. Com k = h$\nu$(1-cos$\theta$)/(m\_e c$^2$), vale k(2+k)/(2+2k+k$^2$) = {[}(1+k)$^2$-1{]}/{[}(1+k)$^2$+1{]}, e a forma explícita da resolução é essa mesma fração multiplicada por (m\_e c$^2$)$^2$ no numerador e no denominador. Testei em Python com valores aleatórios de k e com valores físicos: a diferença entre as formas é nula até o erro de arredondamento. Também conferi que $\surd$((1+$\beta$)/(1-$\beta$)) = 1+k = $\nu$/$\nu$'. A resolução chega ao resultado pela transformação de Lorentz do quadrivetor (E, p), com $\beta$ com sinal, que dá $\beta$\_u < 0. O gabarito usa direto a fórmula do Doppler relativístico para aproximação. As duas abordagens são válidas.%
\end{itemize}%
\paragraph{Pontos de dúvida}%
\begin{itemize}%
\item%
\textit{(key)} Na resolução detalhada do gabarito, os parênteses estão errados na passagem final: aparece '(1 + h$\nu$/(m\_ec\^{}2)(1{-}cos$\theta$)\^{}2 {-} 1)', quando o correto é '(1 + h$\nu$(1-cos$\theta$)/(m\_e c$^2$))$^2$ - 1', com o quadrado aplicado a (1+k) inteiro. O mesmo vale para o denominador. A resposta final do gabarito, listada por item, está correta, então o problema é só de formatação.%
\item%
\textit{(key)} Também na lei dos cossenos do gabarito aparece 'p\_e\^{}2 = h$\nu$/c\^{}2 + (h$\nu$')/c\^{}2 - ...', quando o correto é (h$\nu$/c)$^2$ + (h$\nu$'/c)$^2$ - ... A conta seguinte, multiplicada por c$^2$, está certa, então é só erro de digitação.%
\item%
\textit{(statement)} O enunciado diz apenas que o observador se move 'na mesma direção do fóton espalhado', sem dizer o sentido. Como $\nu$' < $\nu$, a única solução física é o observador ir ao encontro do fóton (sentido oposto ao de propagação). A resolução e o gabarito interpretam assim e respondem com V em módulo e o sentido indicado. Um aluno que tomasse 'mesma direção' como 'mesmo sentido' não acharia solução real, mas isso não afeta a correção desta resolução.%
\end{itemize}%
\end{document}